\documentclass[11pt]{article}
\usepackage{fontspec,amsmath,amssymb,amsthm,hyperref}
\usepackage[margin=1in]{geometry}
\setlength{\emergencystretch}{3em}
\title{Transfer the Gibbs score estimate to an L\ensuremath{{}^{2}} reference}
\author{LangevinLean mathematical exposition}
\date{}
\begin{document}\maketitle
\noindent Source module: \texttt{UnadjustedLangevin/StationaryMoments.lean}. Compile this proof document with XeLaTeX or LuaLaTeX. The accompanying HTML pairs the mathematical argument with the exact Lean statements and records the scope of the formalization.


A reference variable \(Y\) whose law is \(\pi\) inherits \(\|\nabla U(Y)\|_{L^2}\le\sqrt{Ld}\). The assertion that its law is Gibbs is a displayed premise here.


\section{lpDrift norm le of gibbs law}


The representative of \texttt{lpDrift} equals \(\nabla U(Y)\) almost everywhere. Change variables using \(P\circ Y^{-1}=\pi_U\), so its squared \(L^2(P)\) norm is \(\int\|\nabla U\|^2d\pi_U\le Ld\). Nonnegativity of the norm and of \(Ld\) then gives \(\|\nabla U(Y)\|_2\le\sqrt{Ld}\). The law identity, rather than independence, is the input to this transfer.



\noindent Lean declarations: \texttt{lpDrift\_norm\_le\_of\_gibbs\_law}.

\section{lpDrift norm le of strongConvex gibbs law}


Use strong convexity to prove \(Z_U<\infty\), and apply the previous law-transfer theorem with that proof. The conclusion is the same \(\sqrt{Ld}\) bound for the lifted drift of any \(L^2\) variable with Gibbs law. There is no separate assumed drift-moment bound hidden in this specialization.



\noindent Lean declarations: \texttt{lpDrift\_norm\_le\_of\_strongConvex\_gibbs\_law}.

\end{document}
